\chapter{Hyperthermia}

\section*{Definition}
Hyperthermia is a core body temperature \ensuremath{>}38.5\textdegree{}C (101.3\textdegree{}F) resulting from failed thermoregulation, where heat production or environmental heat overwhelms the body's heat dissipation capacity. It is distinct from fever, as the hypothalamic set point remains normal and fever-reducing medications are ineffective.

\section*{What Happens in Your Body}
Heat accumulates when thermal load (exogenous from environment or endogenous from metabolism/exercise) exceeds heat dissipation through radiation, conduction, convection, and evaporation. Hyperthermia causes direct cellular damage through protein denaturation, membrane disruption, and thermolabile enzyme inactivation, triggering a systemic inflammatory response with endothelial injury, coagulopathy, and multi-organ failure. Heat stroke involves core temperature \ensuremath{>}40\textdegree{}C (104\textdegree{}F) with CNS dysfunction.

\section*{Causes}
\begin{itemize}
  \item \textbf{Environmental/classic:} Heat wave, high humidity, lack of air conditioning, urban heat island effect, inadequate fluid intake
  \item \textbf{Exertional:} Strenuous exercise in hot or humid conditions (athletes, military personnel, laborers)
  \item \textbf{Drug-induced:} Anticholinergics (impaired sweating), sympathomimetics (cocaine, MDMA), neuroleptic malignant syndrome (antipsychotics), serotonin syndrome (SSRIs, MAOIs), malignant hyperthermia (volatile anesthetics, succinylcholine)
  \item \textbf{Endocrine:} Thyroid storm, pheochromocytoma
  \item \textbf{Neurologic:} Hypothalamic stroke or lesion, status epilepticus, CNS infection
  \item \textbf{Risk factors:} Extremes of age, obesity, cardiovascular disease, dehydration, alcohol, anhidrosis, impaired mobility
\end{itemize}

\section*{When to See a Doctor}
See your doctor if:
\begin{itemize}
  \item Core temperature \ensuremath{>}40\textdegree{}C (104\textdegree{}F) with altered mental status (heat stroke) --- medical emergency
  \item Hot, dry skin with anhidrosis (classic heat stroke) or profuse sweating (exertional heat stroke)
  \item Seizure, obtundation, or coma
  \item Persistent vomiting, hypotension, or signs of multi-organ dysfunction (AKI, hepatic injury, DIC, rhabdomyolysis)
\end{itemize}

\section*{Related Symptoms}
\begin{itemize}
  \item Heat cramps, heat exhaustion, heat stroke, rhabdomyolysis, acute kidney injury, DIC
\end{itemize}
\textbf{Important:} Heat stroke mortality ranges from 10--65\% and is directly correlated with the duration and degree of temperature elevation; immediate whole-body cooling is the priority before any laboratory or imaging evaluation.

\section*{Self-Care / Home Management}
\begin{itemize}
  \item Move to a cool, shaded, or air-conditioned environment immediately
  \item Remove excess clothing and apply cool water spray with fanning for evaporative cooling
  \item Hydrate with cool water or electrolyte drinks if conscious and not vomiting
\end{itemize}

\section*{How Your Doctor Will Evaluate You}
\begin{itemize}
  \item Rectal core temperature measurement (tympanic, oral, and axillary unreliable in hyperthermia)
  \item Complete metabolic panel, CK, troponin, coagulation profile (PT/PTT, INR, fibrinogen, D-dimer), CBC, lactate, blood cultures
  \item Urinalysis for myoglobin (rhabdomyolysis); ECG for dysrhythmias; CT head if CNS pathology suspected
\end{itemize}

\section*{Treatment Options}
\begin{itemize}
  \item Immediate cooling: Cold water immersion (gold standard for exertional heat stroke), evaporative cooling (mist \& fan), ice packs to neck/axillae/groin, cold IV saline, or intravascular cooling catheters
  \item Goal: Core temperature \ensuremath{<}38.9\textdegree{}C (102\textdegree{}F) within 30 minutes
  \item Airway management, IV fluids, and vasopressors as needed; monitor for multi-organ dysfunction
  \item Fever-reducing medications (acetaminophen, NSAIDs) are ineffective and potentially harmful in hyperthermia
  \item Dantrolene: Specific treatment for malignant hyperthermia; benzodiazepines for neuroleptic malignant syndrome and serotonin syndrome
\end{itemize}

\section*{References}

\begin{thebibliography}{99}
\bibitem{harrison-hypertherm} Longo DL, et al. \emph{Harrison's Principles of Internal Medicine}. 21st ed. McGraw-Hill; 2022. Chapter 476: Hyperthermia and Hypothermia.
\bibitem{uptodate-hypertherm} Bouchama A, et al. Heat stroke and hyperthermia. In: UpToDate, Post TW (Ed), UpToDate, Waltham, MA; 2025.
\bibitem{bouchama} Bouchama A, et al. Heat stroke: a clinical review. \emph{N Engl J Med}. 2022;386(21):2001-2012.
\bibitem{cialdini} Cialdini K, et al. Malignant hyperthermia: diagnosis and management. \emph{Anesth Analg}. 2023;136(5):876-889.
\bibitem{rosen} Walls RM, et al. \emph{Rosen's Emergency Medicine: Concepts and Clinical Practice}. 10th ed. Elsevier; 2023. Chapter 135: Heat-Related Illness.
\bibitem{lefeber} Lefeber EJ, et al. Neuroleptic malignant syndrome: a clinical update. \emph{Am J Psychiatry}. 2023;180(6):424-435.
\end{thebibliography}
